% Part: proof-theory
% Chapter: normalization
% Section: permutations

\documentclass[../../../include/open-logic-section]{subfiles}

\begin{document}

\olfileid{pt}{nor}{per}

\olsection{বিন্যাস-রূপান্তর}

\Log{N1i}-তে কর্তন হলো এমন খণ্ড
যার শেষ সূত্রসংঘটনটি
!!a{elimination} বিধির প্রধান
পূর্বধারণা। স্বাভাবিকীকরণের প্রমাণে
একটি ক্ষেত্র হলো সর্বাধিক কর্তনগুলির
দৈর্ঘ্য কমানো; তাতে !!{proof}-এর
কর্তনদৈর্ঘ্যও কমে। দৈর্ঘ্য~$>1$-এর
কর্তন নানাভাবে শেষ হতে পারে: কর্তনে
$!A$-এর শেষ !!{formula}-সংঘটনটি
\Elim{\lor} বা \Elim{\lexists}-এর
সিদ্ধান্ত কি না এবং সেটি কোন
!!{elimination} বিধির প্রধান
পূর্বধারণা, তার উপর তা নির্ভর করে।

যেমন, সংশ্লিষ্ট বিধি যদি \Elim{\lor}
ও \Elim{\land} হয়, $!A \ident !B \land !C$
হয় এবং একটি উপ-!!{proof}~$\delta_1$
নিচের মতো শেষ হয়:
\[
\AxiomC{}
\RightLabel{$\delta_2$}
\DeduceC{$!D \lor !E$}
\AxiomC{$\Discharge{!D}{x}$}
\RightLabel{$\delta_3$}
\DeduceC{$!B \land !C$}
\AxiomC{$\Discharge{!E}{y}$}
\RightLabel{$\delta_4$}
\DeduceC{$!B \land !C$}
\DischargeRule{\Elim{\lor}}{x\,y}
\TrinaryInfC{$!B \land !C$}
\RightLabel{\Elim{\land}[1]}
\UnaryInfC{$!B$}
\DisplayProof
\]
% BN-SRC-577 TARGET-CORRECTION-BEGIN
\par\smallskip\noindent\textnormal{\small\textbf{উৎস-সংশোধন (BN-SRC-577):} উদাহরণের কর্তনসূত্র উৎসগদ্যে $!A\ident !B\lor !C$ হলেও প্রদর্শিত \Elim{\lor}-সিদ্ধান্ত ও পরের \Elim{\land}-এর প্রধান পূর্বধারণা $!B\land !C$। তাই এই ক্ষেত্রের $!A$-কে $!B\land !C$ করা হয়েছে।}
% BN-SRC-577 TARGET-CORRECTION-END
কর্তনে শেষ !!{formula}-সংঘটন~$!A_n$
হলো \Elim{\lor}-এর সিদ্ধান্ত—নিচের
$!B \land !C$। শেষের আগের
!!{formula}-সংঘটন~$!A_{n-1}$ হলো
\Elim{\lor}-এর গৌণ পূর্বধারণাগুলির
একটি।

এই কর্তনের দৈর্ঘ্য কমাতে
\Elim{\lor} ও \Elim{\land}
বিধির ক্রম বদলাতে, অর্থাৎ
বিন্যাস করতে পারি। তখন
$\delta$-তে উপপ্রমাণ~$\delta_1$-র
জায়গায় নিচেরটি বসাই:
\[
\AxiomC{}
\RightLabel{$\delta_2$}
\DeduceC{$!D \lor !E$}
\AxiomC{$\Discharge{!D}{x}$}
\RightLabel{$\delta_3$}
\DeduceC{$!B \land !C$}
\RightLabel{\Elim{\land}[1]}
\UnaryInfC{$!B$}
\AxiomC{$\Discharge{!E}{y}$}
\RightLabel{$\delta_4$}
\DeduceC{$!B \land !C$}
\RightLabel{\Elim{\land}[1]}
\UnaryInfC{$!B$}
\DischargeRule{\Elim{\lor}}{x\,y}
\TrinaryInfC{$!B$}
\DisplayProof
\]
\Elim{\lor}-এর নিচে \Elim{\land}-এর
পূর্বধারণায় যে কর্তনগুলি শেষ হতো,
এখন সেগুলি \Elim{\lor}-এর গৌণ
পূর্বধারণাগুলির উপরে থাকা
\Elim{\land} বিধিগুলির কোনো একটির
পূর্বধারণায় শেষ হয়। অর্থাৎ
প্রতিটির দৈর্ঘ্য~$1$ কমেছে।

উপরন্তু, $\delta_2$, $\delta_3$
বা $\delta_4$-র সম্পূর্ণ ভিতরে,
অথবা $\delta$-তে $\delta_1$-র
সম্পূর্ণ বাইরে থাকা কর্তন
অপরিবর্তিত: তাতে আগের একই
!!{formula}s ও একই বিধি থাকে।
তাই $\delta$-র সংশ্লিষ্ট কর্তনের
মাত্রা ও দৈর্ঘ্যও একই থাকে।
সুতরাং নতুন !!{proof}-এর
কর্তনমাত্রা বাড়েনি এবং
নতুন !!{proof}-এর
কর্তনদৈর্ঘ্য পুরোনো
!!{proof}-এর চেয়ে কমেছে।

\begin{prob}
বিন্যাস-রূপান্তরের পরে অন্তত একটি
কর্তনের দৈর্ঘ্য~$1$ কমে।
\Elim{\land}-কে \Elim{\lor}-এর
উপরে সরালে আসলে দুটি কর্তনখণ্ডের
দৈর্ঘ্য কমে; তাই এ ক্ষেত্রে
!!{proof}-এর কর্তনদৈর্ঘ্য অন্তত~$2$
কমে। এমন একটি বিন্যাসেই
!!a{proof}-এর কর্তনদৈর্ঘ্য
$2$-এর বেশি কমতে পারে কি?
কর্তনের \emph{মাত্রা} কমতে পারে কি?
\end{prob}

\Elim{\lor} বা \Elim{\lexists}
বিধির উপরে !!a{elimination}
বিধি সরানোর প্রক্রিয়াকে
\emph{বিন্যাস-রূপান্তর} বলে।
কয়েকটি বিধিজোড়ার ছক
\olref{tab:permutation-conversions}-এ
আছে। (অন্যগুলি অনুশীলনের জন্য
রাখা হয়েছে।)

% \begin{table}
% \begin{multline*}
% \shoveleft{\AxiomC{$!D \lor !E$}
% \AxiomC{$\Discharge{!D}{x}$}
% \shortDeduce
% \DeduceC{$!B \land !C$}
% \AxiomC{$\Discharge{!E}{y}$}
% \shortDeduce
% \DeduceC{$!B \land !C$}
% \DischargeRule{\Elim{\lor}}{x\,y}
% \TrinaryInfC{$!B \land !C$}
% \RightLabel{\Elim{\land}[1]}
% \UnaryInfC{$!B$}
% \DisplayProof}\\
% \shoveright{\longrightarrow\ 
% \AxiomC{$!D \lor !E$}
% \AxiomC{$\Discharge{!D}{x}$}
% \shortDeduce
% \DeduceC{$!B \land !C$}
% \RightLabel{\Elim{\land}[1]}
% \UnaryInfC{$!B$}
% \AxiomC{$\Discharge{!E}{y}$}
% \shortDeduce
% \DeduceC{$!B \land !C$}
% \RightLabel{\Elim{\land}[1]}
% \UnaryInfC{$!B$}
% \DischargeRule{\Elim{\lor}}{x\,y}
% \TrinaryInfC{$!B$}
% \DisplayProof}
% \\
% \shoveleft{\AxiomC{$!D \lor !E$}
% \AxiomC{$\Discharge{!D}{x}$}
% \shortDeduce
% \DeduceC{$!B \lif !C$}
% \AxiomC{$\Discharge{!E}{y}$}
% \shortDeduce
% \DeduceC{$!B \lif !C$}
% \DischargeRule{\Elim{\lor}}{x\,y}
% \TrinaryInfC{$!B \lif !C$}
% \AxiomC{$!B$}
% \RightLabel{\Elim{\lif}}
% \BinaryInfC{$!C$}
% \DisplayProof}\\
% \shoveright{\longrightarrow\ 
% \AxiomC{$!D \lor !E$}
% \AxiomC{$\Discharge{!D}{x}$}
% \shortDeduce
% \DeduceC{$!B \lif !C$}
% \AxiomC{$!B$}
% \RightLabel{\Elim{\lif}}
% \BinaryInfC{$!C$}
% \AxiomC{$\Discharge{!E}{y}$}
% \shortDeduce
% \DeduceC{$!B \lif !C$}
% \AxiomC{$!B$}
% \RightLabel{\Elim{\lif}}
% \BinaryInfC{$!C$}
% \DischargeRule{\Elim{\lor}}{x\,y}
% \TrinaryInfC{$!C$}
% \DisplayProof}
% \\
% \shoveleft{\AxiomC{$!D \lor !E$}
% \AxiomC{$\Discharge{!D}{x}$}
% \shortDeduce
% \DeduceC{$!B \lif !C$}
% \AxiomC{$\Discharge{!E}{y}$}
% \shortDeduce
% \DeduceC{$!B \lif !C$}
% \DischargeRule{\Elim{\lor}}{x\,y}
% \TrinaryInfC{$!B \lif !C$}
% \AxiomC{$!B$}
% \RightLabel{\Elim{\lif}}
% \BinaryInfC{$!C$}
% \DisplayProof}\\
% \shoveright{\longrightarrow\ 
% \AxiomC{$!D \lor !E$}
% \AxiomC{$\Discharge{!D}{x}$}
% \shortDeduce
% \DeduceC{$!B \lif !C$}
% \AxiomC{$!B$}
% \RightLabel{\Elim{\lif}}
% \BinaryInfC{$!C$}
% \AxiomC{$\Discharge{!E}{y}$}
% \shortDeduce
% \DeduceC{$!B \lif !C$}
% \AxiomC{$!B$}
% \RightLabel{\Elim{\lif}}
% \BinaryInfC{$!C$}
% \DischargeRule{\Elim{\lor}}{x\,y}
% \TrinaryInfC{$!C$}
% \DisplayProof}
% \\
% \shoveleft{\AxiomC{$!D \lor !E$}
% \AxiomC{$\Discharge{!D}{x}$}
% \shortDeduce
% \DeduceC{$\lforall[x][!B(x)]$}
% \AxiomC{$\Discharge{!E}{y}$}
% \shortDeduce
% \DeduceC{$\lforall[x][!B(x)]$}
% \DischargeRule{\Elim{\lor}}{x\,y}
% \TrinaryInfC{$\lforall[x][!B(x)]$}
% \RightLabel{\Elim{\lforall}}
% \UnaryInfC{$!B(t)$}
% \DisplayProof}\\
% \shoveright{\longrightarrow\ 
% \AxiomC{$!D \lor !E$}
% \AxiomC{$\Discharge{!D}{x}$}
% \shortDeduce
% \DeduceC{$\lforall[x][!B(x)]$}
% \RightLabel{\Elim{\lforall}}
% \UnaryInfC{$!B(t)$}
% \AxiomC{$\Discharge{!E}{y}$}
% \shortDeduce
% \DeduceC{$\lforall[x][!B(x)]$}
% \RightLabel{\Elim{\lforall}}
% \UnaryInfC{$!B(t)$}
% \DischargeRule{\Elim{\lor}}{x\,y}
% \TrinaryInfC{$!B(t)$}
% \DisplayProof}
% \end{multline*}
% \caption{\Log{N1i}-এর কয়েকটি বিন্যাস-রূপান্তর।}
% \ollabel{tab:permutation-conversions}
% \end{table}

{\small
\begin{longtable}{@{}l@{}}
\AxiomC{$!D \lor !E$}
\AxiomC{$\Discharge{!D}{x}$}
\shortDeduce
\DeduceC{$!B \land !C$}
\AxiomC{$\Discharge{!E}{y}$}
\shortDeduce
\DeduceC{$!B \land !C$}
\DischargeRule{\Elim{\lor}}{x\,y}
\TrinaryInfC{$!B \land !C$}
\RightLabel{\Elim{\land}[1]}
\UnaryInfC{$!B$}
\DisplayProof
$\longrightarrow$\\[6ex]
\quad\AxiomC{$!D \lor !E$}
\AxiomC{$\Discharge{!D}{x}$}
\shortDeduce
\DeduceC{$!B \land !C$}
\RightLabel{\Elim{\land}[1]}
\UnaryInfC{$!B$}
\AxiomC{$\Discharge{!E}{y}$}
\shortDeduce
\DeduceC{$!B \land !C$}
\RightLabel{\Elim{\land}[1]}
\UnaryInfC{$!B$}
\DischargeRule{\Elim{\lor}}{x\,y}
\TrinaryInfC{$!B$}
\DisplayProof
\\[10ex]
\AxiomC{$!D \lor !E$}
\AxiomC{$\Discharge{!D}{x}$}
\shortDeduce
\DeduceC{$!B \lif !C$}
\AxiomC{$\Discharge{!E}{y}$}
\shortDeduce
\DeduceC{$!B \lif !C$}
\DischargeRule{\Elim{\lor}}{x\,y}
\TrinaryInfC{$!B \lif !C$}
\AxiomC{$!B$}
\RightLabel{\Elim{\lif}}
\BinaryInfC{$!C$}
\DisplayProof
$\longrightarrow$\\[6ex]
\AxiomC{$!D \lor !E$}
\AxiomC{$\Discharge{!D}{x}$}
\shortDeduce
\DeduceC{$!B \lif !C$}
\AxiomC{$!B$}
\RightLabel{\Elim{\lif}}
\BinaryInfC{$!C$}
\AxiomC{$\Discharge{!E}{y}$}
\shortDeduce
\DeduceC{$!B \lif !C$}
\AxiomC{$!B$}
\RightLabel{\Elim{\lif}}
\BinaryInfC{$!C$}
\DischargeRule{\Elim{\lor}}{x\,y}
\TrinaryInfC{$!C$}
\DisplayProof
\\[10ex]
\AxiomC{$!D \lor !E$}
\AxiomC{$\Discharge{!D}{x}$}
\shortDeduce
\DeduceC{$!B \lif !C$}
\AxiomC{$\Discharge{!E}{y}$}
\shortDeduce
\DeduceC{$!B \lif !C$}
\DischargeRule{\Elim{\lor}}{x\,y}
\TrinaryInfC{$!B \lif !C$}
\AxiomC{$!B$}
\RightLabel{\Elim{\lif}}
\BinaryInfC{$!C$}
\DisplayProof
$\longrightarrow$\\[6ex]
\AxiomC{$!D \lor !E$}
\AxiomC{$\Discharge{!D}{x}$}
\shortDeduce
\DeduceC{$!B \lif !C$}
\AxiomC{$!B$}
\RightLabel{\Elim{\lif}}
\BinaryInfC{$!C$}
\AxiomC{$\Discharge{!E}{y}$}
\shortDeduce
\DeduceC{$!B \lif !C$}
\AxiomC{$!B$}
\RightLabel{\Elim{\lif}}
\BinaryInfC{$!C$}
\DischargeRule{\Elim{\lor}}{x\,y}
\TrinaryInfC{$!C$}
\DisplayProof
\\[10ex]
\AxiomC{$!D \lor !E$}
\AxiomC{$\Discharge{!D}{x}$}
\shortDeduce
\DeduceC{$\lforall[x][!B(x)]$}
\AxiomC{$\Discharge{!E}{y}$}
\shortDeduce
\DeduceC{$\lforall[x][!B(x)]$}
\DischargeRule{\Elim{\lor}}{x\,y}
\TrinaryInfC{$\lforall[x][!B(x)]$}
\RightLabel{\Elim{\lforall}}
\UnaryInfC{$!B(t)$}
\DisplayProof
\quad$\longrightarrow$\\[8ex]
\quad\AxiomC{$!D \lor !E$}
\AxiomC{$\Discharge{!D}{x}$}
\shortDeduce
\DeduceC{$\lforall[x][!B(x)]$}
\RightLabel{\Elim{\lforall}}
\UnaryInfC{$!B(t)$}
\AxiomC{$\Discharge{!E}{y}$}
\shortDeduce
\DeduceC{$\lforall[x][!B(x)]$}
\RightLabel{\Elim{\lforall}}
\UnaryInfC{$!B(t)$}
\DischargeRule{\Elim{\lor}}{x\,y}
\TrinaryInfC{$!B(t)$}
\DisplayProof
\\
\caption{\Log{N1i}-এর কয়েকটি বিন্যাস-রূপান্তর।}
\ollabel{tab:permutation-conversions}
\end{longtable}
}

\begin{prob}
\Elim{\lor}-এর পরে \Elim{\lor}
বা \Elim{\lnot} থাকলে সংশ্লিষ্ট
বিন্যাস-রূপান্তরগুলি দিন।
\end{prob}

\begin{prob}
\Elim{\lexists}-এর পরে
\Elim{\lexists} থাকলে
বিন্যাস-রূপান্তরগুলি দিন।
এই শেষ ক্ষেত্রে কেন কোনো
আইগেনচল-শর্ত লঙ্ঘিত হয় না,
ব্যাখ্যা করুন।
\end{prob}

বিন্যাস-রূপান্তর প্রয়োগে
আইগেনচল-শর্ত ভাঙছে না,
তা নিশ্চিত হওয়া দরকার।
ধরা যাক,
\[
\AxiomC{}
\RightLabel{$\delta_2$}
\DeduceC{$!D \lor !E$}
\AxiomC{$\Discharge{!D}{x}$}
\RightLabel{$\delta_3$}
\DeduceC{$\lexists[x][!B(x)]$}
\AxiomC{$\Discharge{!E}{y}$}
\RightLabel{$\delta_4$}
\DeduceC{$\lexists[x][!B(x)]$}
\DischargeRule{\Elim{\lor}}{x\,y}
\TrinaryInfC{$\lexists[x][!B(x)]$}
\AxiomC{$\Discharge{!B(c)}{z}$}
\RightLabel{$\delta_5$}
\DeduceC{$!C$}
\DischargeRule{\Elim{\lexists}}{z}
\BinaryInfC{$!C$}
\DisplayProof
\]
এর বদলে নিচেরটি বসালে
\[
\AxiomC{}
\RightLabel{$\delta_2$}
\DeduceC{$!D \lor !E$}
\AxiomC{$\Discharge{!D}{x}$}
\RightLabel{$\delta_3$}
\DeduceC{$\lexists[x][!B(x)]$}
\AxiomC{$\Discharge{!B(c)}{z}$}
\RightLabel{$\delta_5$}
\DeduceC{$!C$}
\DischargeRule{\Elim{\lexists}}{z}
\BinaryInfC{$!C$}
\AxiomC{$\Discharge{!E}{y}$}
\RightLabel{$\delta_4$}
\DeduceC{$\lexists[x][!B(x)]$}
\AxiomC{$\Discharge{!B(c)}{z}$}
\RightLabel{$\delta_5$}
\DeduceC{$!C$}
\DischargeRule{\Elim{\lexists}}{z}
\BinaryInfC{$!C$}
\DischargeRule{\Elim{\lor}}{x\,y}
\TrinaryInfC{$!C$}
\DisplayProof
\]
আশঙ্কা হতে পারে, আইগেনচল~$c$
হয়তো~$!D$-তেই আছে। কারণ মূল
$\Elim{\lexists}$-এর উপরে
অনুমান~$!D$ অবমুক্ত হয়; তাই
মূল !!{proof}-এ
\Elim{\lexists} বিধির প্রধান
পূর্বধারণার উপরে খোলা কোনো
অনুমানে আইগেনচলটি ঘটে না,
এবং সেখানে আইগেনচল-শর্ত
পূরণ হয়। কিন্তু নতুন
!!{proof}-এ \Elim{\lexists}
বিধিগুলির পর পর্যন্ত~$!D$
অবমুক্ত হয় না; তাই শর্তটি
ভাঙার কথা। তবে মনে রাখি,
আমরা !!{proof}s-কে নিয়মিত
ধরেছি: প্রত্যেক আইগেনচল
মাত্র একটি বিধির আইগেনচল
এবং !!{proof}-এ কেবল
\Elim{\lexists}-এর গৌণ
পূর্বধারণার উপরে ঘটে।
বিশেষত~$c$, $\delta_3$
বা~$\delta_4$-তে ঘটতে পারে না।
% BN-SRC-576 TARGET-CORRECTION-BEGIN
\par\smallskip\noindent\textnormal{\small\textbf{উৎস-সংশোধন (BN-SRC-576):} এই N1i অধ্যায়ে অস্তিত্বসূচক নিষ্কাশনের বিধি \Elim{\lexists}; উৎসের অনুশীলন ও আইগেনচল-যুক্তিতে দুটি স্থানে সাধারণ \Elim{\exists} নির্দেশ আছে। প্রদর্শিত বৃক্ষ ও N1i বিধি-সারণির সঙ্গে মিলিয়ে দুটিই \Elim{\lexists} করা হয়েছে।}
% BN-SRC-576 TARGET-CORRECTION-END

এই উদাহরণে আরও একটি বিষয়
দেখা যায়: নতুন !!{proof}-এ
$\delta_5$-এর দুটি অনুলিপি
আছে। $\delta_5$-এ সর্বাধিক
কর্তন থাকলে তার মানে
কর্তনদৈর্ঘ্য আমরা
\emph{কমাইনি}! তাই
$\delta_5$-এ সর্বাধিক কর্তন
নেই জেনেই বিন্যাস-রূপান্তর
প্রয়োগ করতে হবে। এ জন্য
সব সময় \emph{সবচেয়ে উপরের}
সর্বাধিক কর্তন বেছে নিই:
অর্থাৎ যার উপ-!!{proof}-এ
(এখানে~$\delta_1$)
ওই উপ-!!{proof}-এর শেষ
বিধির উপরে শেষ হয়
এমন অন্য সর্বাধিক কর্তন
নেই। ফলে~$\delta_5$-এ
(এমনকি $\delta_2$, $\delta_3$
বা~$\delta_4$-তেও) অন্য
সর্বাধিক কর্তন থাকার
সম্ভাবনা বাদ যায়।

\end{document}
